\documentclass[12pt,a4paper]{report}
\usepackage[utf8]{inputenc}
\usepackage[margin=1in]{geometry}
\usepackage{hyperref}
\usepackage{longtable}
\usepackage{tabularx}
\usepackage{booktabs}
\usepackage{xcolor}

\title{DarkNetTracker Project Report}
\author{Project Maintenance Update}
\date{May 17, 2026}

\begin{document}

\maketitle
\tableofcontents
\newpage

\chapter{Executive Summary}

DarkNetTracker is a research-oriented traffic-correlation platform with a Python analysis engine, an Express backend, SQLite persistence, and a Next.js dashboard. This update focused on closing the highest-priority production gaps that were identified in the project backlog.

The most important outcome of this pass is that the repository now has a realistic production-hardening baseline instead of documentation that overclaimed completion. Security controls, environment-aware startup behavior, migration support, CI validation, and frontend resilience have been improved and documented.

\chapter{Repository Status}

\section{Completed in This Update}

\begin{itemize}
  \item Added a central backend configuration module with development and production settings.
  \item Added a migration runner and schema migration tracking table for SQLite startup.
  \item Added production startup enforcement that blocks demo passwords in production mode.
  \item Replaced ad hoc security headers with \texttt{helmet}.
  \item Tightened CORS to configured origins only.
  \item Added CSRF token generation and verification for authenticated write routes.
  \item Hardened upload validation for size and extension checks.
  \item Added structured request IDs, centralized backend error handling, and graceful shutdown.
  \item Added Python subprocess retry, timeout, and failure logging improvements.
  \item Added Pydantic-based runtime validation for Python pipeline inputs.
  \item Added frontend error and loading boundaries plus accessibility improvements for critical controls.
  \item Added GitHub Actions CI and refreshed project documentation.
  \item Added backend API smoke tests and frontend unit test scaffolding with passing test runs.
  \item Added an admin metrics endpoint and an OpenAPI source endpoint for operations and integration work.
\end{itemize}

\section{Current Status Snapshot}

\begin{longtable}{p{4.2cm}p{3cm}p{6.3cm}}
\toprule
Area & Status & Notes \\
\midrule
Backend API & Improved & Security middleware, validation, migration runner, graceful shutdown \\
Frontend Dashboard & Improved & Error boundary, loading state, auth flow updated for CSRF \\
Python Engine & Improved & Pydantic validation and structured progress logging \\
Documentation & Improved & README, runbook, troubleshooting, architecture, security checklist \\
CI/CD & Added baseline & GitHub Actions workflow for installs, migration, tests, frontend build \\
Testing & Improved & Python validation tests, backend API smoke tests, frontend unit tests, frontend build \\
\bottomrule
\end{longtable}

\chapter{Architecture Update}

\section{Runtime Flow}

\begin{enumerate}
  \item The frontend calls the backend bootstrap endpoint for environment metadata.
  \item A user authenticates and receives a bearer token plus a CSRF token.
  \item Authenticated analysis requests are validated and persisted as sessions and jobs.
  \item The queue launches the Python engine through \texttt{engine\_api.py}.
  \item Results are written to SQLite and broadcast to the frontend over WebSocket.
  \item Export endpoints produce CSV, HTML, and PDF outputs from stored session results.
\end{enumerate}

\section{Security Posture}

\begin{longtable}{p{4.5cm}p{3cm}p{5.5cm}}
\toprule
Control & State & Detail \\
\midrule
Default credentials & Improved & Production startup now fails if demo passwords remain configured \\
CORS restrictions & Implemented & Requests are limited to configured origins \\
CSRF protection & Implemented & Authenticated non-GET routes require \texttt{X-CSRF-Token} \\
Security headers & Implemented & \texttt{helmet} is enabled with CSP and related headers \\
Rate limiting & Implemented & General API, auth, and upload limiters are configured \\
Upload validation & Implemented & File size and extension checks added before processing \\
TLS support & Available & Optional HTTPS startup via configured certificate and key paths \\
\bottomrule
\end{longtable}

\chapter{Implementation Details}

\section{Backend}

The backend now uses a single configuration source in \texttt{backend/config.js}. Startup creates required directories, reads environment settings, and validates production requirements before listening for traffic.

Database initialization was moved behind a migration runner in \texttt{backend/migrations.js}. This gives the project a maintainable path for future schema changes instead of one-off \texttt{ALTER TABLE} calls embedded directly in runtime setup.

The Express server now includes:

\begin{itemize}
  \item \texttt{helmet}-based security headers
  \item request IDs
  \item stricter request validation with Zod
  \item CSRF enforcement on authenticated write operations
  \item upload rate limiting and file validation
  \item centralized error responses
  \item graceful shutdown for the HTTP server and WebSocket clients
  \item an admin metrics endpoint at \texttt{/api/metrics}
  \item an OpenAPI source endpoint at \texttt{/api/docs/openapi.yaml}
\end{itemize}

\section{Frontend}

The Next.js dashboard was updated to support the new CSRF flow by storing the CSRF token returned at login and including it automatically on non-GET API requests. The frontend also now has:

\begin{itemize}
  \item an application error boundary
  \item an application loading screen
  \item accessibility labels for important login and admin controls
  \item baseline unit tests using Vitest and Testing Library
\end{itemize}

\section{Python Engine}

Runtime validation in \texttt{validation.py} now uses a Pydantic model for request shape and boundary enforcement. The pipeline service logs major processing milestones so long runs are easier to trace when launched by the backend queue.

\chapter{Verification}

\section{Checks Run}

\begin{itemize}
  \item Frontend production build completed successfully.
  \item Backend migration bootstrap completed successfully.
  \item Backend startup smoke test completed successfully.
  \item Login, authenticated read, and CSRF-protected logout smoke tests completed successfully.
  \item Python validation test suite completed successfully: 11 tests passed.
  \item Backend API smoke tests completed successfully: 2 tests passed.
  \item Frontend unit tests completed successfully: 3 tests passed.
  \item Python simulate-mode pipeline smoke test completed successfully.
\end{itemize}

\section{Known Remaining Gaps}

\begin{itemize}
  \item Broader automated backend API tests are still needed.
  \item Browser-level end-to-end tests are still needed.
  \item External monitoring, alerting, and dependency vulnerability cleanup are still pending.
  \item The broader feature roadmap remains open and should not be described as complete.
\end{itemize}

\chapter{Conclusion}

This update moves DarkNetTracker from a mostly working prototype with optimistic documentation toward a more credible production-ready baseline. The codebase now has stronger security defaults, a repeatable database migration path, more reliable subprocess execution, cleaner operational documentation, and a frontend that understands the hardened auth flow.

The project is in a significantly better state for continued development and controlled deployment, but it should still be described as an actively maturing platform rather than a fully finished enterprise system.

\end{document}
